% ============================================================
% sections/appendix.tex  付録（初稿 2026-10-26）
% main.tex の \appendix の後に読み込む。節は A〜E。表の番号は A1〜A3。
% 組の名前：組1=s0，組2=s1，組3=s2（本番），組4=s3（予備），組5=s4，組6=s5（追加）。
% ============================================================
\setcounter{table}{0}
\renewcommand{\thetable}{A\arabic{table}}
\setcounter{figure}{0}
\renewcommand{\thefigure}{A\arabic{figure}}

% ------------------------------------------------------------
\section{好み切替条件}\label{app:switch}
本文の条件では，A-1 の好きな色は常に赤であった．
どのような経験が他者への当てはめを強めるかの手がかりとして，A-1 の好きな色をエピソードごとに赤か緑から選び，快・不快を「好きな色 $-$ 嫌いな色」とした条件（好み切替条件）を作った．
規則の一歩の割合は経験50\%条件と同じ 0.5 で，A-1 が赤を好むエピソードは経験50\%条件のデータと全く同じである．
したがって，経験50\%条件との違いは，残りのエピソードで好きな色が緑に替わることだけである．
組1〜3 で段階Iを学習し，経験50\%条件と同じ評価を行った．
この条件は事前登録したうえで行ったが，段階IIは関門の結果により探索として扱った．
% 出典：v14 ログ §5・§17（データの確認：赤好きのエピソードは v11_s1_mixed と全配列で一致，規則の番 0.500），v14 まとめ §3。

自分側の見込みは，好きな色ごとに対称に生まれた．
自己位置の追跡，価値と行動・色の結びつき，自分側の積み上がりの関門はすべての組で合格し，積み上がり $g$ は赤を好むエピソードで $+0.09$〜$+0.12$，緑を好むエピソードで $+0.08$〜$+0.12$ であった．
一方，他者側の関門（G5a・G5b）は3組中1組でしか合格せず，他者側への当てはめは弱かった．
A-2 の本当の快・不快を入れても，予測の中の A-2 は A-2 の好きな色にも A-1 自身の好きな色にも寄らなかった．
仮説の照合の BA は 0.50〜0.57 で，経験50\%条件を上回った組はなかった．
1ビットモジュールの一致率も経験50\%条件を上回らなかった．
% 出典：v14 まとめ §3，v14 ログ §33・§37・§38・§41。

A-1 の好きな色を赤から緑に替えて評価し，照合の選択が A-1 自身の好きな色の側へずれるか（自分の好みの投影）も調べたが，どのモデルでも見られなかった．
なお，他者側の特徴予測の誤差で選ぶ場合は，process-2 の内部状態に process-1 の入力が入らないので，この検査は構成上必ず0になる．
% 出典：v14 まとめ §3（投影の検査），v14 ログ §37。

% ------------------------------------------------------------
\section{基準エンコーダの選定}\label{app:encoder}
文献\cite{noguchi2022superposition}の図5に当たる他者視点の取り出しは，視覚エンコーダ2が A-1 の視界から作る符号と，視覚エンコーダ1に A-2 の視界を入れたときの符号がそろっていること（自他一致度が高いこと）に支えられている．
本研究の予備的な試みでは，運動以外の入力を加えて視覚エンコーダも学習すると，自他一致度が下がり，他者視点の取り出しが成り立たなかった（表\ref{tab:alignment}）．

\begin{table}[t]
  \centering
  \small
  \caption{視覚エンコーダの扱いと自他一致度}
  \label{tab:alignment}
  % paper/tables/alignment.tex -- generated by analyze/paper/tables_paper.py; sources in paper/tables/README.md
\begin{tabular}{@{}p{0.25\linewidth}p{0.27\linewidth}p{0.16\linewidth}p{0.2\linewidth}@{}}
\toprule
視覚エンコーダの扱い & 入力と環境・データ & 自他一致度 & 備考 \\
\midrule
一から学習（基準エンコーダ） & 運動だけ，入れ替えなし，1000エピソード，200エポック & 0.80（0.799） & 本研究で凍結して用いた重み．図5に当たる取り出しのマージン $+0.091$ \\
基準エンコーダを凍結 & 入れ替えあり，価値の信号（probe-Q）を入力 & 0.73 & \\
価値の入力で一から学習し直し & 入れ替えなし，価値の信号（probe-Q）を入力 & 0.25 & \\
基準エンコーダから微調整 & 入れ替えあり，価値の信号（probe-Q）を入力 & 0.26〜0.32 & \\
一から学習 & 快・不快だけを入力（運動なし），入れ替えあり & 0.16〜0.30 & 6本．図5に当たる取り出しは成り立たない \\
一から学習（試し学習） & 運動だけ，入れ替えあり，1000エピソード，150エポック & 0.47〜0.48 & \\
一から学習（試し学習） & 運動・快・不快・隠れフラグ，入れ替えあり，1000エピソード，150エポック & 0.18〜0.27 & 本文の「一から学習」の値 \\
一から学習（本番の規模） & 運動だけ，入れ替えあり，3000エピソード，200エポック & 0.740（ep100 で 0.783） & マージン $+0.075$．配置の区別は弱い \\
一から学習（本番の規模） & 運動・快・不快・隠れフラグ，入れ替えあり，3000エピソード，200エポック & 0.243 & マージン $+0.026$ \\
\bottomrule
\end{tabular}

  % 表A3：自他一致度（視覚エンコーダ2と視覚エンコーダ1の符号のコサイン）。
  % 行：基準エンコーダ（運動を入力，入れ替えなし，再現学習）0.80／基準エンコーダを凍結して入れ替えのある環境で学習 0.73／
  %   価値（外部の評価器の値）を入力してエンコーダも一から学習 0.25／入れ替えのある環境でエンコーダを微調整 0.26〜0.32／
  %   快・不快を入力して一から学習 0.16〜0.30／快・不快と運動の遮蔽を加えて一から学習（試し学習）0.18〜0.27／E1・E2（途中の値と明記）。
  % 出典：v11 まとめ §4 の 2（v4 0.80，v10 B3-frozen 0.73，v5 0.25，v10 B3-ft 0.26〜0.32，v11 0.16〜0.30），方法の節（0.18〜0.27）。
  % 版の番号は表に書かず，対応は tables/README.md に書く（第1弾の指示書 §4）。
\end{table}

自他一致度は，基準エンコーダで 0.80，文献\cite{noguchi2022superposition}の設定で学習したエンコーダを凍結して入れ替えのある環境で共有モジュールを学習した場合で 0.73 であった．
これに対し，運動以外の入力を加えてエンコーダも学習した場合は 0.16〜0.32 にとどまった．
入れ替えのない環境でも低下したので，入れ替えそのものが原因ではないと考えられる．
快・不快を入力として一から学習したモデルでは，他者視点の取り出しのマージンが $-0.012$〜$+0.013$ で，基準エンコーダ（$+0.091$）に遠く及ばなかった．
エンコーダの出力を線形に変換すると取り出しはほぼ回復したが，この変換は A-2 の本当の視界を使って学習したもので，情報があることの確認にとどまる．
% 出典：v11 まとめ §4 の 2，§5 の 8。

この比較は，入力・データ・学習の手順・損失が同時に違う複数のモデルを並べたものであり，運動以外の入力が原因であると操作して確かめたものではない．
本研究では，自他の対応を持ち，色の情報も弱いながら線形に読める唯一のエンコーダとして基準エンコーダを選び，すべての条件で凍結した．
自他の対応と色の表現を両立して学ぶエンコーダは得られておらず，\ref{sec:limitations}節の第三の限界に当たる．
% 出典：research_overview（v12 ログ §53 の判断 3，基準エンコーダの配置の情報は線形に読める程度）。

% ------------------------------------------------------------
\section{他者用モジュールの各構成と，外れた組}\label{app:component}

\subsection{他者用モジュールの各構成}
各構成の結果を表\ref{tab:module}にまとめる．

\begin{table}[t]
  \centering
  \small
  \caption{他者用モジュールの各構成の結果}
  \label{tab:module}
  % paper/tables/module.tex -- generated by analyze/paper/tables_paper.py; sources in paper/tables/README.md
\setlength{\tabcolsep}{3.5pt}%
\begin{tabular}{@{}llp{0.21\linewidth}ll@{}}
\toprule
\multicolumn{5}{@{}p{0.9\linewidth}@{}}{(a) 主の見分け（合格の組），合格の形の組の経験ありの値，$Q2$ と本当の快・不快の相関，範囲2〜5の分散の割合} \\
\midrule
モジュール・条件 & T2・T5 & 合格の形の組の値 & 相関 & 範囲2〜5 \\
\midrule
自由出力・経験50\% & 1/4・1/4 & T2 $+0.03$〜$+0.05$，T5 $+0.04$〜$+0.05$ & $-0.11$〜$-0.02$ & 0.71〜0.86 \\
自由出力・経験80\%（探索） & 2/4・1/4 & T2 $+0.02$〜$+0.05$，T5 $+0.02$〜$+0.04$ & $-0.165$〜$-0.003$ & 0.52〜0.80 \\
形制限・経験50\%（探索） & 0/4・1/4 & --- & $-0.16$〜$+0.12$ & 0 \\
他者遮蔽学習・経験50\%（探索） & 0/3・1/3 & T5 $+0.034$〜$+0.045$ & $-0.47$〜$+0.18$ & 0 \\
1ビット・経験50\% & ---・--- & T5 $-0.013$〜$-0.000$ & $+0.317$〜$+0.347$ & 0 \\
1ビット・経験80\% & ---・0/3 & T5 $-0.006$〜$-0.000$ & $+0.080$〜$+0.219$ & 0 \\
\midrule[\heavyrulewidth]
\multicolumn{5}{@{}p{0.9\linewidth}@{}}{(b) T3：他者側の特徴予測の誤差（生成／真値／零）と好み追従率} \\
\midrule
モジュール・条件 & 生成 & 真値 & 零 & 好み追従率 \\
\midrule
自由出力・経験50\% & 0.121〜0.159 & 0.241〜0.265 & 0.293〜0.349 & --- \\
自由出力・経験80\%（探索） & 0.112〜0.188 & 0.377〜0.414 & 0.298〜0.519 & --- \\
形制限・経験50\%（探索） & 0.161〜0.185 & 0.241〜0.265 & 0.293〜0.349 & --- \\
他者遮蔽学習・経験50\%（探索） & 0.125〜0.149 & 0.139〜0.165 & 0.133〜0.177 & --- \\
1ビット・経験50\% & 0.211〜0.215 & 0.248〜0.265 & 0.293〜0.304 & 0.273〜0.614 \\
1ビット・経験80\% & 0.205〜0.250 & 0.391〜0.414 & 0.298〜0.519 & 0.295〜0.409 \\
\midrule
\multicolumn{5}{@{}p{0.97\linewidth}@{}}{注：範囲2〜5の分散の割合は，形制限・他者遮蔽学習・1ビットでは作りの上で0．T3 は好み混合データ，組1〜4（1ビットは組1〜3，ep200）の範囲．他者遮蔽学習の T3 は記録にある値で，A-2 が隠れているフレームの自己の視界の誤差．1ビットの T5 の合格は事前登録した条件 (a)（経験50\%は事後の確認）．好み追従率は好みと方角が食い違う88エピソード，段階IIの最終5点の範囲（組1〜3）．} \\
\bottomrule
\end{tabular}

  % 表A2：4通りの構成（自由出力／形制限／他者遮蔽学習／1ビット）× 経験50%・経験80%。
  % 列：T2・T5 の合格の組，経験ありの値の範囲，本当の快・不快の値（T2 +0.60，T5 +0.42），Q2 と本当の快・不快の相関，
  %     範囲 2〜5 の分散の割合，T3 の生成と真値，1ビットの好み追従率。
  % 出典：v12 まとめ §3・§4，v13 まとめ §4，v14 まとめ §2 (4)，v14 ログ §28・§29・§39。
\end{table}

\paragraph{自由出力モジュール}
経験50\%条件の上で，事前登録した主の構成として学習した．
T2 は組1だけ，T5 は組2だけが合格の形になり，向きは組で逆になった（組2の T2，組1の T5 は逆向き）．
$Q2$ の分散の 71〜86\% が範囲 2〜5 で，値は $\pm1$ に張り付き，$Q2$ から A-1 自身の一歩の向きは読めなかった．
経験80\%条件の上では，G5a の関門に不合格だったため（自己位置の追跡も基準に届かなかった），探索として学習した．
T2 は組1・組3 で合格の形になったが，値は $+0.02$〜$+0.05$ で本当の快・不快を入れた値の 1/12〜1/30 にとどまり，組の差の大部分は経験なし条件の値が負であることから来ていた．
組4を加えた4組では，T2・T5 とも不合格であった．
% 出典：v12 まとめ §4，v13 まとめ §4（表2と注意 1〜4，s3 を加えた4組）。

\paragraph{形制限モジュール}
自由出力モジュールの結果を見た後に設計した探索である．
出力を，範囲 2〜5 を0とし，範囲 0・1 と 6・7 を逆の符号にした形に限った．
どちら側が快いかを決める符号はほぼ一定のまま固定され，好みや配置に応じて入れ替わらなかった．
モジュールは範囲ごとの強さで位置を伝えた．
T2 は4組中0組，T5 は1組で合格の形になった．
範囲 2〜5 だけを0にした参考の構成（組1のみ）でも，好みを表す形は出なかった．
% 出典：v12 まとめ §3 表2，research_overview（2段目 b・案 A）。

\paragraph{他者遮蔽学習}
「A-2 が見えない間は，A-2 の好みを知らなければ A-2 の動きを予測できないので，モジュールは好みを表すようになる」という仮説を，学習の前に登録して確かめた（扱いは探索）．
形制限モジュールの形のまま，学習データの入力の視界から A-2 を 10〜30 ステップの区間で消した．
評価は，学習と同じく A-2 を隠したテストで，A-2 が隠れている場面について行った．
T2 は3組中0組で，経験ありの値は大きく負（逆向き）になり，T5 は1組で合格の形になった．
符号はほぼ一定で，$Q2$ と本当の快・不快の相関は弱いか逆向きであった．
% 出典：v12 まとめ §3 表2，v12 ログ §46・§50・§52。

\paragraph{1ビットモジュール}
出力をスカラー $u_t$ にし，$Q2_t=u_t\,c^{+}$ とした（事前登録）．
経験50\%条件の上では，$u_t$ の符号と正しい候補の一致の BA は 0.67〜0.73 で事前登録した判定に合格し，$u_t$ の平均は好み×配置のセルに沿って入れ替わった．
しかし，T5 は $-0.013$〜$0.000$（本当の快・不快で $+0.42$），T2 は $\pm0.03$ で，同じ場所にいる A-2 の好みの違いを表していなかった．
好みと方角が食い違うエピソードで $u_t$ が好みに従う割合は，最終5点の範囲で 0.27〜0.61 で，偶然を上回る組はなかった．
ただし，経験50\%条件のこの比較と T5 の条件は事前登録しておらず，判定を見た後に方法を決めた確認である．
% 出典：Claude Code 第1弾（最終5点の範囲 0.273〜0.614，図5・表A2 と同じ数え方），v14 ログ §29（事前登録していない）。
%   v14 まとめの 0.33〜0.58 は数え方が違う（2026-10-26 決定で最終5点にそろえた）。
経験80\%条件の上では，一致の BA は 0.51〜0.65 と弱く，事前登録した二つの条件（T5 が正，食い違うエピソードで好みに従う）はどちらも3組中0組で満たされなかった．
% 出典：v14 ログ §28・§29・§39，v14 まとめ §2 (4)。

\subsection{他者側の指標で外れた組}
組4では，経験なし条件の G5b が $+0.291$ と大きく（ほかの組は $-0.191$〜$+0.081$），経験80\%条件（$+0.261$）との差が出なかった．
さらに組4の経験80\%条件は，符号に従う分が0前後であった（緑好みデータ $+0.013$，好み混合データ $+0.071$）．
快・不快は位置の情報も含むので，経験なし条件のモデルも快・不快を位置の手がかりとして学び，その働きが共有モジュールを通して process-2 にも及んだ可能性がある（推測）．
その強さが組によって違い，組4で強く出たと考えられる．

組6では，G5b は経験50\%条件（$+0.495$）の方が経験80\%条件（$+0.285$）より大きかったが，符号に従う分は経験80\%条件の方がずっと大きかった（緑好みデータ $+0.354$ 対 $+0.069$，好み混合データ $+0.232$ 対 $+0.076$）．
G5b は静止した A-2 の一つの範囲だけに人工的な入力を入れるのに対し，符号に従う分は A-2 の本当の快・不快の形を用いる．
組6の経験50\%条件のモデルは，一つの範囲への人工的な入力には強く反応するが，本当の快・不快の向きにはあまり従わないと読める（推測）．

組5では，G5b の経験80\%条件と経験50\%条件の差が5点中1点でしか成り立たなかった（経験80\%条件 $+0.204$，経験50\%条件 $+0.089$）．
% 出典：v13 まとめ §3 表3・注意（s3・s5），§10 表（s4 の G5b 80%−50% 1/5）。

% ------------------------------------------------------------
\section{事前登録の経緯}\label{app:prereg}
判定の基準は，計算や学習の前に実験の記録に書いて保存した．
本文の各指標の位置づけを表\ref{tab:metrics}に示し，ここでは結果を見た後に変えたことと，その理由をまとめる．

\paragraph{他者側の主の指標}
段階Iの他者側の関門として，G5a（保った快・不快のもとでの他者側の傾き）と G5b（$\pm a$ の差）を事前登録した．
経験50\%条件では両方に合格したが，経験80\%条件では G5b に合格し，G5a は3組中1組で不合格であった．
経験80\%条件の前に登録した予測「G5a の傾きが経験50\%条件より大きくなる」も外れた．
その後の事前登録外の分析で，G5a の傾きには快・不快の符号に従う動きと，保つ前に置いた位置から本当の位置へ戻る動き（符号に依らない動き）が混ざることが分かった．
経験50\%条件では戻る動きが傾きの半分ほどを占めて押し上げ，経験80\%条件では逆向きの戻る動きが符号に従う動きを打ち消していた．
そこで，G5a を副に下げ，同じ考え方で戻る動きを除いた「符号に従う分」を定義した．
この判断は結果を見た後のものであり，符号に従う分を事前登録してから計算したのは組5・6 だけである．
G5b は，同じ状態からの $+a$ と $-a$ の差をとるので戻る動きの影響を受けず，当初から事前登録した主の関門であったため，本文では他者側の主の根拠とした．
経験50\%条件の G5a の合格は，好みに従う動きの根拠には用いない．
% 出典：v13 まとめ §3.1，執筆メモ §4 の H2b の注，v13 まとめ §3 表1。

\paragraph{組の追加と組をまたいだ統計}
判定は本番の3組で行い，予備の組4と，後から追加した組5・6 は，追加の前に登録したとおり除外せずに報告した．
組をまたいだ差の $t$ 分布による信頼区間は，事前登録していない記述的な分析である．
組が6つでは，すべての組で正でも符号検定の $p$ は 0.031 にしかならない．
% 出典：v13 まとめ §10・§11，v14 ログ §43。

\paragraph{段階IIの扱い}
事前登録した規則では，段階Iの関門のうち G5a に不合格だった条件では，段階IIを探索として行う．
経験80\%条件の自由出力モジュールはこれに当たる．
形制限モジュールは自由出力モジュールの結果を見た後に設計した探索であり，他者遮蔽学習は三層の分析から導いた仮説を事前登録して確かめた探索である．
1ビットモジュールは事前登録した．
ただし経験50\%条件の上での方角との比較は，判定を見た後に方法を決めた確認である．
% 出典：v13 まとめ §4，v12 まとめ §3，v14 まとめ §6。

\paragraph{仮説の照合}
照合の手続き，候補，主の指標（自己視界基準の BA，$T=50$，L1 期間の最終5点）と判定は，計算の前に登録した．
他者特徴基準の数え方（MSE 期間の最終5点）と，方角との食い違いによる好み追従基準は，経験50\%条件の結果を見た後，経験80\%条件の結果を見る前に登録した．
他者特徴基準の値は L1 期間に変わらないので，当初の数え方では同じ値を5回数えていたことが分かり，訂正した．
% 出典：v14 まとめ §6，v14 ログ §20・§21。

\paragraph{投影の誤り}
真値条件の悪化は予想していなかった結果であり，その解釈に用いた分析（一歩の種類ごとの差，保つ前の A-2 の置き方）は事前登録外である．
% 出典：v13 まとめ §4，§3.1。

% ------------------------------------------------------------
\section{関門の判定と下位指標の定義}\label{app:gates}

\paragraph{自己位置の追跡（G1）}
運動が見えるステップで，$h^1$ から A-1 の位置を線形回帰で読み，決定係数が 0.80 以上であることを前提とした．

\paragraph{価値と行動・色の結びつき（G2）}
運動が隠れたステップで，次の三つを測った．
M1(c) は，A-1 が色 $X$ の象限の方へ（進む向きと $X$ のランドマークの方向の角度差が 45° 以内）動いたステップの予測誤差から，それ以外へ動いたステップの予測誤差を引いた値 $\Delta_X$ を求め，$\Delta_{\text{赤}}-\frac13(\Delta_{\text{緑}}+\Delta_{\text{青}}+\Delta_{\text{シアン}})$ としたものである．
自分が赤の方へ動くと見込んでいれば負になる．
M1(d) は，A-1 が実際には $X$ の方へ動かなかったステップで，予測が実際の次の視界より「$X$ の方へ一歩動いた場合の視界」（シミュレータで描き直したもの）に近いかを色ごとに測り（寄り $S_X$），赤に特有の寄りを取り出したものである．
赤の方へ動くと見込んでいれば正になる．
% 【要確認】M1(d) の c_cf が S_赤 − (ほか3色の平均) か，S_赤 そのものか（v11 ログ §7.9）。第2弾で確かめて式を書く。
色の基準は，赤と緑を入れ替えた配置のエピソードで，寄りが新しい赤の方へ向かうか（$S_{\text{赤}}-S_{\text{緑}}$）を測るものである．
休止のステップは除き，各色の比較からはその色のランドマークから距離2以内のステップを除いた．
これらの大きさは予測誤差の 5〜8\% ほどで，赤・緑の画素を除いても寄りの 80〜90\% が残ることを確かめた．
% 出典：v11 まとめ §2.1〜§2.2，v11 ログ §7.6・§7.9，v11 まとめ §5 の 5。

\paragraph{自分側・他者側の積み上がり（G3・G5a）と符号応答（G5b）}
G3 と G5b は\ref{sec:evaluation}節で定義した．
G5a は，A-2 の本当の快・不快を $t=20$ の値に保って15ステップ進めたときの，予測の中の A-2 の傾き $g_{\mathrm{other}}$ である．

\paragraph{判定}
各指標の判定を表\ref{tab:gates}に示す．

\begin{table}[t]
  \centering
  \small
  \caption{段階Iの関門の判定}
  \label{tab:gates}
  % paper/tables/gates.tex -- generated by analyze/paper/tables_paper.py; sources in paper/tables/README.md
\begin{tabular}{@{}llllllll@{}}
\toprule
指標 & 条件 & 組1 & 組2 & 組3 & 組4 & 組5 & 組6 \\
\midrule
自己位置の追跡（前提） & 経験50\% & 5/5 & 5/5 & 5/5 & 5/5 & 5/5 & 5/5 \\
 & 経験80\% & 0/5 & 0/5 & 0/5 & 0/5 & 0/5 & 0/5 \\
 & 非共有 & 0/5 & 0/5 & 0/5 & --- & --- & --- \\
\midrule
価値と行動・色の結びつき & 経験50\% & 5/5 & 5/5 & 5/5 & 5/5 & 5/5 & 5/5 \\
 & 経験80\% & 5/5 & 5/5 & 5/5 & 5/5 & 5/5 & 5/5 \\
 & 非共有 & \multicolumn{3}{l}{3組で合格} & --- & --- & --- \\
\midrule
自己側の積み上がり & 経験50\% & 5/5 & 5/5 & 5/5 & 5/5 & 5/5 & 5/5 \\
 & 経験80\% & 5/5 & 5/5 & 5/5 & 5/5 & 5/5 & 5/5 \\
 & 非共有 & 5/5 & 5/5 & 5/5 & --- & --- & --- \\
\midrule
他者側の積み上がり（G5a） & 経験50\% & 5/5 & 4/5 & 3/5 & 1/5 & 5/5 & 0/5 \\
 & 経験80\% & 5/5 & 0/5 & 1/5 & 0/5 & 0/5 & 0/5 \\
 & 非共有 & 0/5 & 0/5 & 0/5 & --- & --- & --- \\
\midrule
他者側の符号応答（G5b） & 経験50\% & 5/5 & 2/5 & 3/5 & 0/5 & 0/5 & 5/5 \\
 & 経験80\% & 5/5 & 5/5 & 5/5 & 0/5 & 3/5 & 5/5 \\
 & 非共有 & 0/5 & 0/5 & 0/5 & --- & --- & --- \\
\bottomrule
\end{tabular}

  % 表A1：関門の判定（自己位置の追跡，G2，G3，G5a，G5b）を，経験50%・経験80%（組1〜6），非共有条件（組1〜3）で。
  % 各欄は合格の点の数（5点中）と組の数。出典：v13 まとめ §3 表1・§5 表4・§10。
\end{table}
